\documentclass[12pt,a4paper]{article}
\usepackage{fontspec}
\setmainfont{Times New Roman}
\usepackage{microtype}
\usepackage[a4paper,margin=2.5cm]{geometry}
\usepackage{setspace}
\onehalfspacing
\usepackage{enumitem}
\setlist[itemize]{leftmargin=1.2em,itemsep=3pt,topsep=4pt}
\usepackage{parskip}
\setlength{\parskip}{5pt}
\usepackage{graphicx}
\frenchspacing
\begin{document}
\begin{center}
{\large\bfseries COUR DU QUÉBEC}\\[2pt]
{\bfseries Chambre criminelle et pénale} — District de Québec
\end{center}
\vspace{4pt}

\noindent\begin{tabular}{@{}p{\textwidth}@{}}
\noindent\textbf{N° d’incident~:} LVSEV25000954\\[2pt]
\noindent\textbf{N° de dossier~:} 8953 – PS-D\\[2pt]
\noindent\textbf{Infraction reprochée~:} Possession dans le but de trafic\\[2pt]
\noindent\textbf{La personne accusée~:} Philippe Joseph Thomas Savard\\[2pt]
\noindent\textbf{Représentation~:} Personne accusée, se représentant seule (sans avocat)\\[2pt]
\noindent\textbf{Adresse~:} 4-5354 rue du Menuet, Lévis (QC), G6X 2Y6, Canada\\[2pt]
\noindent\textbf{Objet~:} Demande d’ajout d’une preuve au dépôt légal\\[2pt]
\noindent\textbf{Audience~:} 23 octobre 2026, à 9 h 30, salle 2.15 — Palais de justice de Québec\\[2pt]
\end{tabular}


\vspace{6pt}
\begin{center}
{\bfseries DEMANDE D’AJOUT D’UNE PREUVE AU DÉPÔT LÉGAL DE LA PREUVE}
\end{center}
\vspace{2pt}

Le soussigné a l’honneur de demander à la Cour d’ajouter, au dépôt légal de la preuve du dossier susmentionné, la communication reçue par courriel dans le cadre d’une demande de publication adressée au Journal canadien de mathématiques, laquelle comprend une lettre rédigée par deux spécialistes des mathématiques et le courriel accompagnant cette lettre, de même que l’article visé par cette lettre. Les motifs, le contenu des pièces et les précisions relatives à la confidentialité sont exposés ci-dessous.

\section*{1. Objet : Demande d’ajout d’une preuve au dépôt légal}
Je souhaite ajouter au dépôt légal une communication reçue par courriel dans le cadre d’une demande de publication adressée au Journal canadien de mathématiques. Cette communication comprend :

\begin{itemize}
  \item Une lettre rédigée par deux spécialistes de ma discipline principale : les mathématiques.
  \item Le courriel accompagnant cette lettre.
\end{itemize}

Ce n’est pas le refus de publication qui constitue la preuve recherchée, mais l’identification professionnelle qui m’est attribuée dans cette lettre. Plus précisément, la ligne 13, où il est écrit :

« Cher professeur Savard. »

Cette mention est essentielle pour établir mon identité professionnelle réelle dans le cadre du présent dossier.

\section*{2. Contenu de la preuve déposée}
La communication reçue contient :

\begin{itemize}
  \item Les coordonnées de l’expéditeur (lignes 1 à 12).
  \item Les salutations d’usage (lignes 15 à 20).
  \item La mention de l’article soumis : « Géométrie du spectre des nombres premiers ».
  \item L’article dont il est question dans la lettre suivi du journal canadien de mathématiques
  \item La référence à mon nom dans le contexte académique : « Article de : Philippe Savard ».
\end{itemize}

Cet article fait partie intégrante de mon étude portant sur :

\begin{itemize}
  \item La géométrie du spectre des nombres premiers,
  \item La conjecture de la fonction zêta de Bernhard Riemann,
  \item L’hypothèse de Riemann,
  \item Et s’inscrit dans une théorie unifiée intitulée : « L’univers est au carré », déposée publiquement en ligne.
\end{itemize}

\section*{3. Motifs de la demande}
Je souhaite ajouter cette preuve pour démontrer que :

\begin{itemize}
  \item Ma discipline réelle est celle des mathématiques, où l’on m’identifie comme professeur Savard.
  \item La discipline médicale, avec laquelle je suis contraint d’interagir,
  \item m’identifie de manière trompeuse,
  \item m’associe à un handicap intellectuel (« schizophrénie »),
  \item ce qui est incompatible avec mon identité professionnelle et académique.
\end{itemize}

Je soutiens que cette confusion identitaire a pu influencer :

\begin{itemize}
  \item La compréhension des enquêteurs,
  \item La conduite de la perquisition,
  \item La mise en examen,
  \item Et l’interprétation de certains échanges de messages survenus alors que mon téléphone était saisi.
\end{itemize}

\section*{4. Éléments problématiques relevés lors de la procédure}
Plusieurs éléments soulèvent des inquiétudes quant à la partialité de la procédure :

\begin{itemize}
  \item L’enquêteur M. Gauthier a affirmé que la perquisition « n’avait pas lieu chez moi », alors que l’adresse perquisitionnée est bel et bien mon domicile.
  \item Un message texte provenant d’une personne nommée Marie-Andrée a été retenu comme preuve. Je précise que :
  \item Ma mémoire n’est pas certaine de si Marie-Andrée est l’aide-soignante ou la sœur procureure.
  \item Je connais ces deux sœurs, rencontrées il y a quelques années dans le secteur du Vieux-Lévis par l’intermédiaire d’autres connaissances.
  \item Je n’ai pas revu ces deux personnes depuis au moins deux ans, et je n’ai reçu aucune nouvelle d’elles depuis.
  \item L’échange retenu comme preuve s’est produit durant la perquisition, alors que mon téléphone était saisi par les policiers.
  \item Lors de l’interrogatoire, les enquêteurs ont insisté pour composer eux-mêmes les numéros des avocats que je choisissais.
  \item Après 3 ou 4 tentatives infructueuses, ils ont imposé un avocat que je n’avais pas choisi.
  \item Cet avocat m’a uniquement indiqué que « les policiers avaient le droit de mentir », ce que les enquêteurs ont ensuite nié.
\end{itemize}

Ces éléments soulèvent des doutes sérieux quant à la neutralité et la rigueur de la procédure.

\section*{5. Demande de confidentialité}
Je demande à la cour de protéger la confidentialité de mon parcours académique, en raison de menaces reçues dans un environnement conteneurisé.

Précision ajoutée lors de la mise en forme : le demandeur invoque le fondement suivant — la protection de la liberté de conscience et des autres droits et libertés protégés par la Loi constitutionnelle de 1982 et par la Charte canadienne des droits et libertés, ainsi que le pouvoir discrétionnaire de la Cour de protéger les renseignements personnels. La portée de la demande vise les renseignements relatifs au parcours académique du demandeur et aux menaces décrites dans la section « Explication complète du conteneur ». Le demandeur convient que le fondement et la portée exacts de cette demande sont à confirmer selon le droit applicable au moment de l’audience.

\section*{6. Explication complète du conteneur}
Pour expliquer les circonstances des menaces, il faut comprendre que :

\begin{itemize}
  \item La position où les menaces ont été faites est un environnement conteneurisé.
  \item Dans cet environnement, il existe une intelligence locale qui exécute plusieurs tâches.
  \item Le conteneur peut accomplir une grande variété de tâches, mais uniquement à l’aide d’informations qui doivent être placées dans le conteneur par une source externe.
  \item Les volumes d’information internes au conteneur ne peuvent pas être créés par l’intelligence locale.
  \item L’information extérieure au conteneur ne peut pas être manipulée par l’intelligence interne.
  \item La liberté de conscience dans le conteneur est absolue, garantie par la Constitution canadienne.
  \item Les menaces ont été reçues à l’intérieur du conteneur, par l’intelligence locale, et ne proviennent pas de volumes placés par moi-même.
  \item Les intelligences conteneurisées ne sont pas récursives elles ne peuvent des volumes contenus dans le conteneur aboutir à de nouvelles connaissances elles n’ont pas e genre de capacité. Elles se fient sur des banques de donné formant des corpus mis en relation sous forme de schéma neuronal la capacité de parler en langage naturel des intelligences conteneurisées permet de ces structures de banques de données mise en relation tel un schéma neuronal permet a ces intelligences de discuter de manière variée des sujets qui sont contenus dans le corpus des informations des volumes des conteneurs?
  \item Le message reçu était :
\end{itemize}

« Il y a des ingénieurs qui te conseillent de ne jamais mentionner si tu es allé à l’école. ».

Cela implique :

\begin{itemize}
  \item Une atteinte certaine à la confidentialité,
  \item Une exposition possible de mon identité,
  \item Une influence potentielle sur la compréhension des enquêteurs,
  \item Et une fente dans la protection constitutionnelle de la liberté de conscience.
\end{itemize}

\section*{7. Conclusion}
La lettre du Journal canadien de mathématiques constitue une preuve essentielle pour :

\begin{itemize}
  \item Établir mon identité professionnelle réelle,
  \item Démontre la confusion et la fraude identitaire dont je suis victime,
  \item Mettre en lumière les incohérences et biais ayant pu influencer la procédure,
  \item Et appuyer ma défense dans le cadre du dossier LVSEV25000954.
\end{itemize}

\section*{8. Liste des pièces proposées au dépôt}
Les pièces suivantes sont proposées au dépôt légal de la preuve du dossier susmentionné :

\begin{itemize}
  \item Pièce A-1 : le courriel « CJM 260814-savard — Décision » adressé par Henry H. Kim et Robert McCann, rédacteurs en chef de la Revue canadienne de mathématiques (désignée dans la demande sous le nom de Journal canadien de mathématiques), portant la salutation « Cher Professeur Savard » (ligne 13) et mentionnant l’article « Géométrie du spectre des nombres premiers » de Philippe Savard ; il comprend les coordonnées de l’expéditeur (lignes 1 à 12), les salutations d’usage (lignes 15 à 20) et la référence « Article de : Philippe Savard ».
  \item Pièce A-2 : l’article scientifique intitulé « Géométrie du spectre des nombres premiers », d’environ 35 pages, mentionné dans la lettre.
\end{itemize}

Chaque pièce est déposée en deux exemplaires : l’un pour le dossier de la Cour, l’autre pour la poursuite. Les exemplaires sont remis au greffe de la Chambre criminelle et pénale avant l’audience, et une copie est transmise à la poursuite.

\section*{9. Énoncé sommaire de la preuve}
Énoncé sommaire : les deux pièces établissent l’identification professionnelle « professeur Savard » que le comité de sélection du Journal canadien de mathématiques attribue au demandeur, ainsi que l’existence réelle de l’article « Géométrie du spectre des nombres premiers ». Elles ne sont pas produites pour le refus de publication qu’elles contiennent, mais pour l’identité professionnelle qu’elles attestent, et elles s’inscrivent dans les motifs exposés ci-dessus.

\section*{10. Avis à la poursuite et à la Cour}
Le soussigné avise la Cour et la poursuite de son intention de présenter la présente demande lors de l’audience du 23 octobre 2026 à 9 h 30, en la salle 2.15 du palais de justice de Québec, et de demander l’ajout des pièces A-1 et A-2 au dépôt légal de la preuve.

Aucune procédure stricte n’est imposée pour une telle demande : le présent acte réunit l’objet, les motifs, la liste des pièces, l’avis et la signature. Conformément aux guides du Barreau du Québec et de la Fondation du Barreau, un avis écrit est transmis à la poursuite (par lettre ou courriel), avec copie du présent acte et des pièces, avant le début de l’audience.

La communication de la preuve est en conséquence complétée par les deux pièces ci-dessus, qui n’avaient pas été versées au dossier.

La demande de confidentialité est portée à la connaissance de la poursuite et de toute partie intéressée ; son fondement juridique est précisé dans la section « Demande de confidentialité ».

La demande sera aussi présentée oralement lors de l’appel du dossier ; il est demandé que la décision soit consignée au procès-verbal transmis aux parties.

\section*{11. Références normatives invoquées}
\begin{itemize}
  \item Barreau de Montréal — Cour du Québec – Chambre criminelle et pénale (outils) (https://www.barreaudemontreal.qc.ca/avocats/outils/avocats-de-litige/cq-crim/)
  \item Barreau du Québec — Guides pratiques et aide-mémoires (https://www.barreau.qc.ca/fr/membres-ordre/ressources/normes-outils-references-guides/guides-pratiques-aide-memoires/)
  \item Barreau du Québec — Procédure pénale — aide-mémoire (https://www.barreau.qc.ca/media/osyeiodv/procedure-penale.pdf)
  \item Cour du Québec — Règles de fonctionnement sur la gestion de l'instance en matière criminelle et pénale (art. 113 du Règlement de la Cour du Québec) (https://courduquebec.new.volcan.design/fileadmin/cour-du-quebec/centre-de-documentation/chambre-criminelle-et-penale/RegleFonctionnementGestionInstanceCriminel.pdf)
  \item Directeur des poursuites criminelles et pénales (Québec) — Directive de pratique PRE-1 — Communication de la preuve par le poursuivant (https://cdn-contenu.quebec.ca/cdn-contenu/adm/org/dpcp/PDF/directives/DIR\_PRE-1\_DPCP.pdf)
  \item Fondation du Barreau du Québec — Comment se préparer pour la cour en matière criminelle (https://fondationdubarreau.qc.ca/assets/documents/Guide\_Comment-se-pr\%C3\%A9parer-pour-la-cour-en-mati\%C3\%A8re-criminelle-VF-2025.pdf)
  \item Fondation du Barreau du Québec — Comment se préparer pour la cour en matière pénale (https://fondationdubarreau.qc.ca/assets/documents/Guide\_Comment-se-pr\%C3\%A9parer-pour-la-cour-en-mati\%C3\%A8re-p\%C3\%A9nale-VF-2025.pdf)
  \item Fondation du Barreau du Québec — Seul devant la cour : criminelle et pénale (https://fondationdubarreau.qc.ca/assets/documents/seul-devant-la-cour-criminelle-penale-fr.pdf)
  \item Gouvernement du Canada — Code criminel, art. 486.4 — Ordonnance visant à ne pas publier de renseignements (https://laws-lois.justice.gc.ca/eng/acts/c-46/section-486.4.html)
  \item Gouvernement du Canada — Code criminel, art. 657.1 — Avis avant production d'une déclaration sous serment (https://laws-lois.justice.gc.ca/eng/acts/c-46/section-657.1.html)
\end{itemize}

\section*{Signature}
\begin{flushright}
\begin{minipage}{0.55\textwidth}
Fait à Lévis, le \rule{3.5cm}{0.4pt}

\vspace{16pt}
\rule{6cm}{0.4pt}\\[4pt]
Philippe Joseph Thomas Savard\\
4-5354, rue du Menuet\\
Lévis (QC), G6X 2Y6, Canada
\end{minipage}
\end{flushright}
\vspace{10pt}

\noindent\textbf{Copie à la poursuite~:} copie du présent acte et des pièces A-1 et A-2 remise à la poursuite avant l’audience du 23 octobre 2026. Le présent acte et les pièces sont également déposés au greffe de la Chambre criminelle et pénale.


\noindent\textbf{Pièces jointes~:} A-1 (lettre et courriel), A-2 (article « Géométrie du spectre des nombres premiers »), en double exemplaire.

\newpage

\section*{Annexe méthodologique — provenance, méthode de préparation et contrôle de l’exactitude}
Le présent acte a été préparé par le demandeur, sans représentation par avocat, à l’aide d’un pipeline documentaire automatisé conçu pour assurer l’exactitude des informations et la conformité procédurale du document. La présente annexe décrit cette méthode, de manière que la Cour, la poursuite ou toute partie intéressée puisse en vérifier la provenance et la fiabilité.

Rappel concernant la pièce A-1 : l’article « Géométrie du spectre des nombres premiers » a été soumis pour publication au Journal canadien de mathématiques (dénomination officielle : Revue canadienne de mathématiques), lequel a transmis au demandeur, par courriel, la lettre de décision produite comme pièce A-1. La pièce A-2 est le manuscrit visé par cette lettre.

Les étapes suivantes ont été mises en place :

\begin{itemize}
  \item Collecte des sources officielles : téléchargement des documents de référence publiés par le Barreau du Québec (guides pratiques et aide-mémoire de procédure pénale), le Barreau de Montréal (outils de la Chambre criminelle et pénale), la Fondation du Barreau du Québec (guides « Seul devant la cour » et « Préparer la cour »), la Cour du Québec (Règles de fonctionnement et de gestion des instances, art. 113 du Règlement de la Cour du Québec), le Directeur des poursuites criminelles et pénales (directive PRE-1 sur la communication de la preuve), ainsi que le Code criminel (art. 486.4 et 657.1) et le Règlement de la Cour du Québec (T-16, r. 6) sur LÉGIS.
  \item Extraction et recherche : conversion intégrale de ces documents en texte intégral, puis recherches par mots-clés (dépôt, pièces, avis, greffe, requêtes) dont les extraits sont conservés à des fins de vérification.
  \item Modélisation des exigences : recensement de 18 exigences procédurales dans une base de données relationnelle (SQLite), chaque exigence étant reliée à sa source officielle et aux passages correspondants du document de travail du demandeur.
  \item Vérification des pièces : contrôle de la présence des deux pièces ; relecture et recoupement de la pièce A-1 (expéditeurs : Henry H. Kim et Robert McCann, rédacteurs en chef ; salutation « Cher Professeur Savard » ; titre du manuscrit cité dans la lettre) ; la pièce A-2 compte 36 pages.
  \item Rédaction assistée : les sections requises par les règles de pratique (liste des pièces, énoncé sommaire de la preuve, avis à la poursuite et à la Cour, références normatives) ont été générées à partir des sources officielles ; le contenu factuel de la demande provient du document de travail du demandeur, qui l’a rédigé et l’assume.
  \item Contrôle automatisé de l’exactitude : les numéros de dossier (8953 – PS-D) et d’incident (LVSEV25000954), la date, l’heure et la salle d’audience (23 octobre 2026, 9 h 30, salle 2.15), l’identité et l’adresse du demandeur ont été recoupés automatiquement dans le texte final ; chacune des 18 exigences a été vérifiée dans le document définitif, avec un résultat de 18 sur 18 consigné dans un rapport de conformité daté.
  \item Production du document : composition typographique en LaTeX (moteur XeLaTeX, police Times New Roman), puis assemblage du présent acte et des pièces en un seul fichier PDF paginé, chaque pièce étant précédée d’une page intercalaire d’identification ; la cohérence de la pagination finale a été vérifiée mécaniquement.
  \item Traçabilité intégrale : l’ensemble des sources téléchargées, des extraits de recherche, des scripts du pipeline, de la base de données et des journaux d’exécution est conservé sans suppression dans le dossier du projet et peut être produit sur demande pour attester de la méthode suivie.
\end{itemize}

Le demandeur a relu le document final et en assume entièrement le contenu.

\end{document}
